\section{ASM and Inline ASM}

In this project, \textbf{Assembly code} and \textbf{Inline Assembly} should be used with extreme caution. While it can offer performance improvements or enable access to low-level hardware features, its use can make code harder to understand, maintain, and port. Always prefer writing in C unless absolutely necessary.

\subsection{General Guidelines for ASM and Inline ASM}
\begin{itemize}
    \item \textbf{Avoid ASM When Possible}: Always try to write in C before resorting to assembly language. C compilers are highly optimized and can often produce efficient machine code that rivals hand-written assembly.
    \item \textbf{Document ASM Code Thoroughly}: Assembly code is often harder to understand than C, so it should be accompanied by thorough comments that explain \textbf{why} ASM was used, as well as what the code does.
\end{itemize}

\begin{lstlisting}[language=bash]
/* Set register to zero using ASM */
mov r0, #0      /* Initialize r0 to 0 */
\end{lstlisting}

\begin{itemize}
    \item \textbf{Use Inline ASM Sparingly}: Inline assembly allows you to embed assembly instructions directly in C code. It should only be used when absolutely necessary, such as when interacting with hardware registers or performing CPU-specific optimizations.
    \item \textbf{Constrain the Scope of Inline ASM}: When using inline assembly, constrain its scope to as few lines as possible and avoid intermixing it with complex C logic. This helps isolate potential issues and reduces the risk of bugs.
\end{itemize}

\begin{lstlisting}[language=C]
void set_flag(void)
{
  asm volatile("mov r0, #1" : : : "r0");  /* Set flag in register r0 */
}
\end{lstlisting}

\begin{itemize}
    \item \textbf{Always Use \texttt{volatile} for Inline ASM}: Inline assembly should be marked as \texttt{volatile} to prevent the compiler from optimizing it out.
    \item \textbf{Do Not Use Inline ASM for Optimizations}: Modern C compilers perform aggressive optimizations. Inline ASM should not be used to try and optimize code performance unless absolutely necessary, as it is often more error-prone than compiler optimizations.
    \item \textbf{Use Constraints Correctly}: When using inline ASM in C, use the appropriate constraints to inform the compiler how the assembly instructions interact with C variables.
\end{itemize}

\begin{lstlisting}[language=C]
int add(int a, int b)
{
  int result;
  asm volatile(
    "add %[r], %[a], %[b]"
    : [r] "=r" (result)    /* Output */
    : [a] "r" (a), [b] "r" (b) /* Inputs */
  );
  return result;
}
\end{lstlisting}

\subsection{When to Use ASM and Inline ASM}
\begin{itemize}
    \item \textbf{Hardware Access}: Use ASM when direct hardware access is required and C alone cannot provide the necessary control (e.g., setting or clearing specific registers).
    \item \textbf{Performance Critical Code}: ASM may be justified in performance-critical sections where C compiler optimizations fall short, but this should be documented and used as a last resort.
    \item \textbf{CPU-Specific Instructions}: Use ASM for executing instructions that are unique to the target CPU (e.g., specific ARM or x86 instructions that are not exposed in C).
\end{itemize}
